\begin{abstract}
Este trabajo aborda el Problema del Clique Máximo (MCP), un problema clásico de optimización combinatoria clasificado como NP-duro, mediante una formulación de Programación Lineal Entera (ILP) resuelta con herramientas de modelado algebraico. Se implementa el modelo matemático en Python utilizando el framework Pyomo \cite{pyomo2021} y el solver de alto rendimiento HiGHS \cite{highs2024}, que internamente aplica el método de Ramificación y Acotamiento (Branch \& Bound) con relajación LP, planos de corte y heurísticas. La metodología se evalúa sobre la instancia \textit{keller4} del benchmark DIMACS \cite{mascia_dimacs}, un grafo de 171 vértices y 9\,435 aristas basado en la conjetura de Keller sobre teselaciones con hipercubos. Los resultados experimentales confirman la obtención de la solución óptima $\omega = 11$ en aproximadamente 12 segundos. El análisis de la relajación LP revela una brecha de integralidad del 87.1\%, característica de la debilidad intrínseca de la formulación LP estándar para el MCP. Se concluye que la combinación de modelado algebraico declarativo con solvers MIP modernos constituye un enfoque efectivo y reproducible para la resolución exacta de problemas de optimización combinatoria.

\vspace{0.4cm}
\noindent \textbf{Palabras clave:} Clique Máximo, Programación Lineal Entera, Branch \& Bound, Pyomo, Optimización Combinatoria, NP-duro, DIMACS.
\end{abstract}